\chapter{Travel Planning and Policies}

\section{When Travel Is Required}
Competition teams (SAE, ASCE, AIAA, Formula SAE, EcoCAR), field survey projects, and client site visits may require travel. All travel must be planned, budgeted, and approved well in advance.

\section{The Go/No-Go Decision}
Before any travel spending, the PA must approve a \textbf{Go/No-Go decision} with the Capstone Leadership Team (CLT). The Go/No-Go considers: is the trip essential? Is the budget available? Is the team prepared? The PA determines the \textbf{travel roster}---the minimum number of travelers needed for the mission.

\section{PA and Student Responsibilities}
\textbf{PA}: confirm whether staff must travel (competition rules), determine roster, review and approve budget, support team logistics, secure OneCard for airfare/rental bookings.

\textbf{Student Travel Lead}: create and manage the travel budget, follow all Mines and Capstone travel policies, serve as primary logistics contact, manage confirmations and check-in details, coordinate flight selection with FROSCH, verify traveler names match government IDs, ensure on-time arrivals.

\section{Budget Planning for Travel}
Prepare a draft budget \emph{before} meeting with CP:
\begin{itemize}[nosep,leftmargin=1.5em]
\item Multiple lodging options with links and cost comparisons
\item Airfare per traveler (include FROSCH service fee if applicable)
\item Ground transport: rental car, mileage for personal vehicles, fuel estimates (note: fuel not reimbursable for personal vehicles)
\item Per diem for all travelers for all days
\item Registration and competition fees
\item Shipping for the project/equipment
\end{itemize}
Use the Budget Template on Canvas. Review the ``Example of a Great Team Budget.''

\section{Key Travel Policies}
\begin{itemize}[nosep,leftmargin=1.5em]
\item \textbf{Carpooling required}. Multiple personal vehicles for the same trip will not be reimbursed.
\item \textbf{Mileage}: Mines rate (\$0.60/mile as of 2/26/24---confirm current rate). Calculate from 1500 Illinois St., Golden CO 80401. Save the Google Maps route as PDF. \textbf{Gas not reimbursable} for personal vehicles.
\item \textbf{Meals}: itemized receipts required. \textbf{No alcohol.} Receipts containing alcohol rejected \emph{in full}---no partial reimbursement.
\item \textbf{Post-travel receipts}: submit within \textbf{one month} of trip completion.
\item \textbf{Airfare}: use FROSCH (university travel agency) when required. Coordinate with Travel Lead.
\end{itemize}

\section{Competition-Specific Considerations}
Registration deadlines (often months ahead), equipment shipping logistics and costs, team shirt production (InnoHub, competition teams only), and the competition schedule's impact on the academic calendar. Plan early; competition travel budgets are among the largest expenditures in capstone.


\section{Travel Budget Template}
A well-organized travel budget includes these categories with line-item detail:

\begin{center}\small
\begin{tabularx}{\textwidth}{lXr}
\toprule
\textbf{Category} & \textbf{Items} & \textbf{Est.\ Cost} \\
\midrule
Airfare & Round-trip per person $\times$ N travelers & \$XXX \\
Lodging & Hotel rate $\times$ nights $\times$ rooms (2 per room) & \$XXX \\
Ground Transport & Rental car + fuel, OR mileage (personal vehicle) & \$XXX \\
Meals & Per diem $\times$ days $\times$ travelers & \$XXX \\
Registration & Competition/event registration fees & \$XXX \\
Shipping & Project equipment shipping (outbound + return) & \$XXX \\
Miscellaneous & Parking, tolls, public transit & \$XXX \\
\midrule
\textbf{Total} & & \textbf{\$XXX} \\
\bottomrule
\end{tabularx}
\end{center}

Get actual quotes (not guesses) for airfare, lodging, and registration. Use university-contracted rates where available. Include links to the specific hotels and flight options you are proposing.

\section{The Post-Trip Checklist}
After returning from travel, complete these tasks within one week:

\begin{enumerate}[nosep,leftmargin=1.5em]
\item Collect all receipts from all travelers. Ensure every receipt is itemized and includes the date, vendor name, and items purchased.
\item Verify no receipts contain alcohol (check carefully---some restaurants list it under cryptic names).
\item Complete the mileage reimbursement form with Google Maps PDF documentation.
\item Complete the reimbursement forms for all out-of-pocket expenses.
\item Submit all forms and receipts to CP within one month.
\item Update the Budget Tracker with actual travel costs.
\item Debrief with the team: what went well, what logistics should improve for next time?
\end{enumerate}

\begin{warningbox}[The Alcohol Receipt Rule Is Absolute]
If \emph{any} receipt includes alcohol---even one drink on a \$150 team dinner receipt---the \textbf{entire receipt is rejected}. There is no partial reimbursement and no exception to this policy. Teams learn this the hard way every semester. The solution: pay for alcohol separately on a personal card, or do not order alcohol on team travel.
\end{warningbox}



\section{Travel Safety and Logistics}
\textbf{Emergency contacts}: before departing, share the team's travel itinerary with your PA and ensure every traveler has the PA's phone number. For competition travel, identify the nearest hospital to the competition venue.

\textbf{Equipment transport}: if shipping the prototype or competition vehicle, plan the logistics well in advance. Measure and weigh the shipment. Get shipping quotes from multiple carriers (FedEx, UPS, freight carriers for large items). Ship early enough to arrive before you do; having your prototype arrive after the competition is a preventable disaster.

\textbf{Personal vehicle travel}: when driving personal vehicles, follow the Mines driving policy (available from EHS). Get adequate sleep before long drives. Switch drivers every 2--3 hours. Do not drive more than 10 hours in a day. Ensure vehicle insurance is current and covers the trip.

\textbf{Hotel logistics}: book rooms with two occupants per room to control costs. Reserve rooms with a free cancellation policy in case plans change. Verify the hotel is within reasonable distance of the competition venue or client site. Request a late checkout on the last day if your event runs late.

\textbf{Documentation during travel}: photograph and video-record everything relevant---competition performance, client site conditions, field survey data, team activities. This content is invaluable for reports, presentations, and the Showcase poster.

\section{Competition-Specific Travel Considerations}
Competition travel has unique requirements beyond standard capstone travel:

\begin{itemize}[nosep,leftmargin=1.5em]
\item \textbf{Registration deadlines}: competition registration often closes months before the event. The Travel Lead must track deadlines and submit registration as soon as the Go/No-Go is approved.
\item \textbf{Tech inspection preparation}: many competitions require pre-inspection paperwork (safety forms, technical declarations, driver certifications). Complete these before traveling.
\item \textbf{Pit area setup}: plan what tools, spare parts, and equipment you need at the competition. Create a packing list and check it before departing.
\item \textbf{Team coordination at the event}: designate roles for the competition day---who operates the prototype, who handles documentation, who interfaces with judges, who manages logistics.
\item \textbf{Post-competition debrief}: before leaving the venue, conduct a hot debrief while memories are fresh. What went well? What failed? What should be improved for next year?
\end{itemize}



\section{International Travel Considerations}
Some competitions and client projects involve international travel. Additional requirements:

\textbf{Passport}: verify all travelers have valid passports with at least 6 months remaining before expiration. Passport processing takes 6--8 weeks (standard) or 2--3 weeks (expedited). Apply early.

\textbf{Mines international travel registration}: Mines requires all university-affiliated international travel to be registered through the Office of Global Education. Complete the registration at least 4 weeks before departure.

\textbf{Travel insurance}: Mines provides institutional travel insurance for registered university travel. Verify coverage details with the Office of Global Education.

\textbf{Equipment export}: if transporting technical equipment or prototypes internationally, verify that no export control restrictions apply. The US Department of Commerce Export Administration Regulations (EAR) and the International Traffic in Arms Regulations (ITAR) restrict export of certain technologies. When in doubt, consult the Mines Office of Research Administration.

\section{What Happens When Plans Change}
Travel plans rarely survive first contact with reality. Common disruptions and responses:

\textbf{Flight cancellation}: rebook through the airline or FROSCH immediately. Document any additional costs (hotel night, meals) for reimbursement. Notify your PA and the competition organizer.

\textbf{Team member illness}: if a team member cannot travel, determine whether the mission can be completed with fewer people. Adjust roles accordingly. The PA may need to approve a modified travel roster.

\textbf{Equipment damage during transport}: this is why you ship the prototype separately and insure it. If damage occurs, assess whether repair is possible on-site. Bring a basic tool kit and spare parts for the most failure-prone components.

\textbf{Weather delays}: build a buffer day into the travel schedule when possible. Arriving the day before the competition starts allows recovery from a weather delay without missing the event.



\section{The Pre-Trip Checklist}
One week before departure, verify:

\begin{itemize}[nosep,leftmargin=1.5em]
\item All travelers have confirmed flight/hotel reservations (names match IDs).
\item Rental car reservation confirmed (driver has valid license and is $\geq$21 years old).
\item Competition registration confirmed and all pre-event paperwork submitted.
\item Prototype packed and shipping confirmed (if shipping separately).
\item Tool kit packed: basic hand tools, multimeter, soldering iron, spare components, zip ties, tape, USB cables.
\item Printed copies of: travel itinerary, hotel confirmation, rental car confirmation, competition schedule, campus map, emergency contacts.
\item All team members have PA phone number and competition emergency contact.
\item Sufficient cash or card for meals and incidentals (receipts will be needed for reimbursement).
\item Camera/phone charged for documentation.
\item Team shirts packed (competition teams).
\end{itemize}

\section{Travel Contingency Planning}
Anticipate problems and have backup plans:

\textbf{If the prototype does not work at the competition}: bring a basic repair toolkit and spare parts for the most failure-prone components. Identify a workspace near the venue (hotel room, parking lot with a table) where repairs can be attempted.

\textbf{If a team member cannot travel}: redistribute roles before departure. Ensure at least two people understand every critical subsystem so the team can function with one member absent.

\textbf{If weather delays travel}: build a buffer day into the schedule. For competitions with fixed start times, this buffer is essential. An extra hotel night (\$50--100) is cheap insurance against missing the event entirely.

\textbf{If costs exceed the budget}: track spending in real time during the trip. If you are trending over budget on meals, switch to grocery store food for the remaining days. If an unexpected cost arises (parking fees, toll roads, equipment replacement), document it immediately for the reimbursement process.



\section{Rental Car and Personal Vehicle Policies}
\textbf{Rental cars}: the PA or a designated student with a OneCard can book a rental car through the university's contracted rental agency. All drivers must be listed on the rental agreement. Most rental companies require drivers to be $\geq$21 years old; some charge a surcharge for drivers under 25. Verify the fuel policy (return full, prepay, etc.) and follow it precisely---rental fuel surcharges are expensive.

\textbf{Personal vehicles}: Mines reimburses mileage at the official rate (confirm with CP before the trip). Mileage is calculated from the Mines campus (1500 Illinois Street, Golden, CO 80401) to the destination. \textbf{Save the Google Maps route as a PDF before departure}---this is required documentation for reimbursement. Gas purchases are \textbf{not reimbursable} for personal vehicles (the mileage rate covers fuel). Carpooling is required; multiple personal vehicles for the same trip will not be reimbursed.

\textbf{Parking and tolls}: reimbursable with receipts. Use hotel parking if available (often included in the room rate). Airport parking receipts must show entry and exit dates matching the travel dates.

\section{Documenting the Trip for Academic and Project Credit}
Travel is not just logistics; it is an opportunity to gather project data and professional experience:

\textbf{Competition teams}: photograph and video-record every event run, tech inspection interaction, and scoring session. Document equipment performance under competition conditions. Record failure modes and repair procedures. This content feeds directly into the FDR report.

\textbf{Client site visits}: photograph the operating environment, existing equipment, user workflows, and any conditions that affect your design requirements. Take notes during every client conversation. These observations often reveal requirements that were not captured during remote meetings.

\textbf{Field survey teams}: follow your data collection protocol precisely. Document GPS coordinates, weather conditions, equipment calibration, and any deviations from the planned procedure. Raw field data that lacks metadata (where, when, how, by whom) cannot be used in analysis.



\section{The Travel Planning Timeline}
Travel planning should begin as soon as the project involves a potential trip:

\textbf{8--12 weeks before}: identify the travel need (competition, client site visit, field survey). Discuss with PA. Initiate the Go/No-Go conversation with CLT.

\textbf{6--8 weeks before}: designate the Travel Lead. Begin budget development. Research lodging and airfare options. Register for the event if applicable.

\textbf{4--6 weeks before}: submit the travel budget to PA and CP for approval. Book lodging. Begin airfare booking (through FROSCH if required). Confirm the travel roster with PA.

\textbf{2--4 weeks before}: finalize all bookings. Prepare the equipment packing list. Ensure all travelers have required documents (ID, registration confirmation, competition paperwork). Ship equipment if required.

\textbf{1 week before}: final pre-trip meeting. Distribute the itinerary to all travelers and PA. Review the pre-trip checklist. Confirm all reservations.

\textbf{During the trip}: document everything. Track spending. Follow all Mines travel policies.

\textbf{Within 1 month after}: submit all receipts and reimbursement forms. Update the Budget Tracker. Debrief with the team.

Teams that wait until 2 weeks before the trip to start planning consistently encounter: sold-out hotels, expensive last-minute airfare, missed registration deadlines, and avoidable logistical crises.

\section{Delegating Travel Duties}
For multi-day trips, distribute logistics across the team to prevent the Travel Lead from being overwhelmed:

\textbf{Travel Lead}: overall coordination, budget management, flight booking.

\textbf{Lodging Coordinator}: hotel research, booking, check-in/check-out.

\textbf{Ground Transport Lead}: rental car booking, driving schedule, fuel stops.

\textbf{Equipment Manager}: packing list, shipping, setup at venue, teardown.

\textbf{Documentation Lead}: photos, videos, competition results, travel receipts collection.

Each role has specific responsibilities, but all roles share the responsibility of being on time, communicating proactively, and following Mines travel policies.



\section{Lessons from Past Capstone Travel}
These lessons come from actual capstone team experiences at Mines:

\textbf{The forgotten adapter}: a team traveled to a competition with their laptop charger but forgot the USB-C cable needed to program the ESP32. They spent 2 hours finding a Best Buy in an unfamiliar city. \emph{Lesson: create a tool/cable checklist specific to your project and verify it the day before departure.}

\textbf{The unviewable video}: a team recorded excellent competition footage on a phone that ran out of storage mid-day. They deleted ``unnecessary'' videos to free space, including the one showing their best performance run. \emph{Lesson: bring a dedicated camera or ensure ample phone storage. Back up footage to a laptop each evening.}

\textbf{The tax receipt}: a team submitted a \$187 dinner receipt for reimbursement. The receipt included two glasses of wine totaling \$24. The entire \$187 was rejected---not just the \$24. The team absorbed the loss personally. \emph{Lesson: the alcohol policy is absolute. Pay for alcohol separately on a personal card, every time.}

\textbf{The missed flight}: a team member's driver's license was expired. TSA would not let them through security. They missed the flight, reboked at their own expense (\$380), and arrived a day late. \emph{Lesson: the Travel Lead must verify that all traveler IDs are valid and match the name on the booking, at least 1 week before departure.}

\textbf{The shipping delay}: a team shipped their prototype via ground shipping to arrive 3 days before the competition. A weather delay added 2 days. The prototype arrived the morning of the competition with no time for setup testing. \emph{Lesson: ship equipment at least 5 business days before the event, or use expedited shipping with tracking.}

\section{The Equipment Packing List}
For competition and field trips, prepare a project-specific packing list:

\textbf{Essential for all trips}: laptop + charger, USB cables for all project hardware, multimeter, basic hand tools (screwdrivers, pliers, wire strippers), soldering iron + solder + flux (if electrical work may be needed), electrical tape + zip ties, spare batteries/power supplies, printed copies of all documentation (schematics, firmware, O\&M procedures), camera for documentation.

\textbf{Competition-specific}: spare components for failure-prone items (fuses, connectors, motors), team shirts, competition rulebook (printed), tech inspection paperwork, safety equipment required by competition rules, display materials (poster, banner, table cover).

\textbf{Field survey-specific}: GPS unit or phone with offline maps, field notebook (waterproof), measurement instruments with calibration certificates, sample containers and labels, weather-appropriate clothing, first aid kit, water and snacks, emergency communication device.



\section{Post-Competition and Post-Trip Documentation}
The work does not end when you return to campus:

\textbf{Financial closeout}: collect all receipts from all travelers within 3 days of return. The CFO organizes receipts by category, verifies no alcohol receipts, completes all reimbursement forms, and submits to CP within 2 weeks (well within the 1-month deadline). Update the Budget Tracker with actual travel costs.

\textbf{Technical debrief}: within 1 week of return, hold a structured team debrief: (1)~what went well (specific successes to replicate), (2)~what failed (specific failures to analyze and prevent), (3)~what we would do differently (actionable improvements for next year's team or for your own future professional practice). Document the debrief in your project repository.

\textbf{Client communication}: if the trip involved a client site visit, send a follow-up email within 48 hours summarizing: observations, data collected, decisions made, and next steps. Attach relevant photos and documentation.

\textbf{Competition results documentation}: for competition teams, document: overall placement, scores in each category, judge feedback (often provided informally---write it down immediately), photos of the prototype in competition, and any equipment issues or damage. This documentation feeds into the FDR report and provides valuable data for future teams.

\textbf{``Lessons for Next Year'' document}: competition teams should create a brief (1--2 page) document capturing specific advice for the next year's team. Include: what to pack, what to order early, which competition rules caused problems, hotel and restaurant recommendations, and any logistical tips specific to the venue. Store this in the Digital Design Archive so future teams can find it.



\section{Per Diem and Meal Management}
Per diem is the daily allowance for meals and incidental expenses during travel. Mines per diem rates follow federal GSA guidelines (rates vary by city):

\textbf{How per diem works}: the per diem rate is a maximum, not an entitlement. You are reimbursed for actual meal costs up to the per diem rate, with itemized receipts. If the per diem for your destination is \$59/day and you spend \$40, you are reimbursed \$40. If you spend \$70, you are reimbursed \$59.

\textbf{Meal receipt requirements}: every meal receipt must be: (1)~itemized (showing each item ordered, not just a total), (2)~free of alcohol, (3)~dated within the travel period, and (4)~from a restaurant or food vendor (grocery store receipts for meal supplies are acceptable). A credit card slip showing only a total amount is not an itemized receipt---request the detailed bill from the server.

\textbf{Group meals}: when the team eats together, one receipt is acceptable with a note listing all team members present. However, the per diem allocation is per person---a \$180 dinner for 4 people at \$45/person is within most per diem rates, but a \$300 dinner for 4 people may exceed the per diem even though the total seems reasonable.

\textbf{The alcohol rule (again)}: this rule is worth repeating because teams violate it every semester. If any item on the receipt is an alcoholic beverage, the \textbf{entire receipt} is rejected for reimbursement. Not just the alcohol---the entire meal. The solution: (1)~do not order alcohol on team travel, or (2)~if you choose to drink, pay for the alcoholic beverages on a separate, personal transaction with a separate receipt. Ask the server for a split check \emph{before} ordering.

\section{Coordinating Large Group Travel}
For teams of 5+ or multi-team travel (e.g., EcoCAR teams sending 8--12 students):

\textbf{Central booking}: designate one person (Travel Lead or PA with OneCard) to make all hotel and flight reservations. This ensures consistent pricing, group rates, and organized documentation. Having 6 people make individual hotel bookings produces 6 different hotels, 6 different rates, and 6 different cancellation policies.

\textbf{Shared transportation}: for driving trips, arrange carpooling with the minimum number of vehicles. Each vehicle needs: a designated driver, a backup driver, a phone with navigation, a charged phone for emergencies, and knowledge of the route. Plan rest stops every 2--3 hours.

\textbf{Communication during travel}: create a group text/chat for real-time coordination. Designate one person as the ``home base'' contact (typically the PA) who has the full itinerary and can coordinate if a vehicle breaks down, a flight is cancelled, or a team member needs assistance.

\textbf{Emergency protocols}: every traveler should have: the PA's phone number, the campus emergency number, the hotel address and phone number, and basic insurance information. For competition travel, add the competition organizer's emergency contact.



\section{Domestic vs. International Travel Differences}
Most capstone travel is domestic (within the US), but some competitions and client projects involve international travel. Key differences:

\textbf{Documentation}: domestic travel requires a valid government-issued ID (driver's license or state ID). International travel requires a valid passport with at least 6 months remaining before expiration. Some countries require visas---check requirements at least 3 months before departure.

\textbf{Institutional requirements}: Mines requires all international travel to be registered through the Office of Global Education. Complete the registration at least 4 weeks before departure. This registration activates institutional travel insurance coverage.

\textbf{Customs considerations}: if transporting a prototype or technical equipment internationally, verify that no export control restrictions apply (ITAR, EAR). Carry a letter from the Capstone Program Director describing the equipment and its educational purpose---this can expedite customs clearance.

\textbf{Communication}: international cell phone service is expensive. Activate an international plan before departure or purchase a local SIM card at your destination. Ensure at least one team member has reliable communication at all times.



\section{Travel as a Professional Development Opportunity}
Travel in capstone is not just logistics; it is a professional development experience that prepares you for business travel throughout your career:

\textbf{Professional behavior while traveling}: you represent Colorado School of Mines and the Capstone Design Program at every interaction during travel---at the airport, the hotel, the competition venue, and the restaurant. Dress professionally when interacting with competition officials, judges, and other teams. Be punctual. Be courteous to service staff. Your behavior reflects on the program and on future Mines teams.

\textbf{Networking at competitions and conferences}: competitions bring together students from dozens of universities. This is a networking opportunity: introduce yourself to other teams, ask about their design approaches, and exchange contact information. The student you meet at an SAE competition may be your colleague at an automotive company in 5 years.

\textbf{Observing other teams}: at competitions, observe how top-performing teams organize their pit area, present to judges, and recover from technical failures. These observations provide lessons that improve your own team's performance. Document what you observe---specific, actionable observations are more valuable than vague impressions.

\textbf{Post-travel reflection}: after returning, reflect (individually and as a team) on: what professional skills did you practice during travel? What would you do differently? How does the competition or site visit experience change your perspective on the project? This reflection feeds into your end-of-semester Individual Performance Evaluation narrative.


\section{Practice Questions}
\begin{enumerate}[leftmargin=1.5em]
\item Create a travel budget for a 4-person team attending a 3-day SAE competition in Detroit. Include all categories with realistic estimates.
\item A dinner receipt includes two beers. What happens to the entire receipt?
\item You forgot to save the Google Maps PDF before a client site visit. How do you document mileage?
\end{enumerate}


